% appendices/A_glossary.tex
\chapter{Beginner Glossary}
\label{app:glossary}



This glossary explains common terms used throughout the manual.

\section*{Core concepts}
\addcontentsline{toc}{section}{Core concepts}

\begin{description}[leftmargin=!,labelwidth=4cm]
  \item[Terminal] A text-based interface where you type commands for the operating system.
  \item[Command] A short instruction you type in the terminal (e.g., \texttt{sudo apt update}).
  \item[sudo] ``Run as administrator'' on Linux. Required for system-level changes.
  \item[Package manager (APT)] A tool that installs and updates software on Debian/Ubuntu (\texttt{apt}).
  \item[Repository] A source of software packages that \texttt{apt} can download from.
  \item[GPG key] A cryptographic key used to verify that packages come from the real publisher and were not modified.
\end{description}

\section*{Docker terms}
\addcontentsline{toc}{section}{Docker terms}

\begin{description}[leftmargin=!,labelwidth=4cm]
  \item[Docker] A system to run applications in containers.
  \item[Container] A lightweight isolated environment that runs an application with its dependencies.
  \item[Image] A packaged template used to create containers (like an app ``installer'', but for containers).
  \item[Volume] A folder where containers store persistent data on the host disk.
  \item[Docker Compose] A tool to define and run multi-container setups using a YAML file.
  \item[Port] A network entry point (e.g., \texttt{8080}) used to reach a service from a browser or an app.
\end{description}

\section*{Remote access and security}
\addcontentsline{toc}{section}{Remote access and security}

\begin{description}[leftmargin=!,labelwidth=4cm]
  \item[SSH] Secure remote login to a Linux machine (terminal access over the network).
  \item[Firewall] A system that controls which network traffic is allowed in/out.
  \item[VPN] A private network tunnel that lets you access services securely as if you were on the local network.
\end{description} 